\chapter[Testes e validação]{Testes e validação} \label{capitulo14}

Os testes do sistema foram planejados com o objetivo de verificar o funcionamento das principais funcionalidades implementadas, a integração entre as camadas da aplicação e a adequação do sistema às necessidades relacionadas ao apoio à avaliação psicopedagógica. A estratégia de validação foi estruturada em duas frentes: testes técnicos realizados durante o desenvolvimento e avaliação de uso com profissionais da área de psicopedagogia.

No backend, foram utilizados Vitest e Supertest para a realização de testes automatizados das funcionalidades da API. Esses testes contemplaram os principais fluxos relacionados à autenticação, gerenciamento de usuários, aprendentes, avaliações, aplicações e relatórios, além da validação dos dados recebidos pela API e do controle de acesso às funcionalidades protegidas.

No frontend, foram utilizados Vitest e React Testing Library, permitindo verificar o comportamento dos componentes e das páginas da aplicação. Foram considerados, entre outros aspectos, os fluxos de autenticação, navegação entre as funcionalidades, utilização de formulários, apresentação dos dados e aplicação das regras de acesso conforme o perfil do usuário.

Além dos testes automatizados, foram realizados testes funcionais e exploratórios durante o desenvolvimento, buscando verificar o funcionamento integrado das funcionalidades do sistema. Nessa etapa, foram observados fluxos como autenticação, cadastro e consulta de aprendentes, criação e gerenciamento de instrumentos de avaliação, realização de aplicações, registro de respostas e geração de relatórios. Também foram verificadas situações relacionadas à consistência dos dados e ao acesso às funcionalidades protegidas.

Como complemento aos testes técnicos, foi planejada uma avaliação com usuários representantes do público-alvo do sistema, especificamente profissionais da área de psicopedagogia. Essa etapa tem como finalidade verificar a percepção dos profissionais quanto à utilização da aplicação e sua adequação às atividades relacionadas à avaliação psicopedagógica.

Para essa avaliação, será disponibilizado um questionário estruturado após a utilização do sistema pelos participantes. O instrumento de coleta buscará reunir informações sobre aspectos como facilidade de uso, organização das informações, clareza das funcionalidades, navegação, adequação das funcionalidades às atividades profissionais e percepção sobre a utilidade do sistema como ferramenta de apoio à avaliação psicopedagógica. Também será disponibilizado espaço para que os participantes registrem dificuldades encontradas, sugestões de melhoria e observações sobre as funcionalidades avaliadas.

A aplicação do questionário permitirá complementar os resultados obtidos nos testes técnicos com a perspectiva dos usuários finais, possibilitando identificar aspectos que não são necessariamente observáveis por meio de testes automatizados, como dificuldades de compreensão, problemas de navegação e adequação do fluxo de trabalho às necessidades dos profissionais.

Os resultados obtidos nessa etapa serão analisados de forma descritiva, considerando as respostas dos participantes e os comentários registrados no questionário. A análise será utilizada para identificar pontos positivos, limitações e oportunidades de melhoria no sistema. Essa avaliação será considerada uma etapa de validação da solução desenvolvida, sem caracterizar uma validação clínica ou metodológica dos instrumentos psicopedagógicos utilizados.

Dessa forma, os testes técnicos verificam o funcionamento e a consistência da implementação, enquanto a avaliação com psicopedagogos acrescenta uma perspectiva de uso e adequação ao contexto profissional. Essa distinção é importante porque o sistema foi desenvolvido como uma ferramenta de apoio à organização e ao registro do processo avaliativo, não tendo como finalidade realizar diagnósticos ou substituir a interpretação realizada pelo profissional.
